\documentclass[zihao=-4,fontset=ubuntu]{graduateart}
\usepackage{amsmath,amsthm,amssymb,amsfonts}

\definecolor{nuanbai}{HTML}{F5F5F5}
\pagecolor{nuanbai!5}
\usepackage{ninecolors}
\usepackage[colorlinks,pdfstartview=FitH,CJKbookmarks=true,bookmarksnumbered=true,bookmarksopen=true,linkcolor=black,filecolor=black,anchorcolor=black,linktocpage=true,urlcolor=black,citecolor=black,pdfencoding={auto},pdfauthor={ Shilong.Lu},pdfsubject={article},pdftitle={titlename},pdfkeywords={article},breaklinks,unicode,linktoc=all,plainpages=false,pdfusetitle,]{hyperref}
\setlength{\headheight}{13.59999pt}
\usepackage{appendix}
\usepackage[backend=biber,style=gb7714-2015,gbpub=false,sorting=nyt,gbbiblabel=bracket,gbcitelabel=circle,gblanorder=chinesehead,citexref=true,gbfootbib=true,gbfootbibfmt=true]{biblatex}
\addbibresource{ref.bib}
\usepackage{booktabs}
\usepackage{multirow}
\usepackage{booktabs}
\usepackage{tabularx}
\usepackage{longtable}
\usepackage{enumitem}
\usepackage{capt-of}


    
\begin{document}
\author{邱~雪~倩}
\teacher{覃~秀~红\quad 副教授、译审}
\subject{泰~语~口~译}
\academy{东~南~亚~语~言~文~化~学~院}
\finishedtime{2025年5月}
\grade{2022级}
\subtitle{在这里填写你的论文题目}
\logo{LOGO.jpg}
\makecover
\input{资料页/undertitlepage}
\input{资料页/titlepage}

\pagenumbering{Roman} % 罗马字母页码
\input{资料页/abstract}
\input{资料页/enabstarct}
\tableofcontents\clearpage % 生成正文目录
\listoffigures\clearpage % 生成图片目录
\listoftables\clearpage % 生成表格目录

\pagenumbering{arabic} % 罗马数字页码
\input{正文页/001}\newpage
\input{正文页/002}\newpage
\input{正文页/003}\newpage
\input{正文页/004}\newpage
\input{正文页/005}\newpage

\newpage


\appendix
\section{附录一}
\input{附录内容/同传材料1 中文}\newpage
\input{附录内容/同传材料1 原文}\newpage
\input{附录内容/同传材料1 修改译文}\newpage
\section{附录二}
\input{附录内容/会议手册}\newpage











\newpage
\begingroup
    \centering
    \printbibliography[heading=bibliography,title=参考文献]
\endgroup


\newpage
\input{附录内容/后记}
\end{document}
